% Build: tectonic d2_decision_model.tex  (XeLaTeX toolchain; re-runs and
% bibtex are handled automatically). Committed PDF is regenerated from this
% source only.
\documentclass[11pt,a4paper]{article}

\usepackage[margin=2.6cm]{geometry}
\usepackage{amsmath}
\usepackage{amssymb}
\usepackage{mathtools}
\usepackage{booktabs}
\usepackage{array}
\usepackage{longtable}
\usepackage{enumitem}
\usepackage[dvipsnames]{xcolor}
\usepackage[colorlinks=true,linkcolor=NavyBlue,citecolor=OliveGreen,urlcolor=blue]{hyperref}
\usepackage[nameinlink,capitalise,noabbrev]{cleveref}
\usepackage{microtype}

\newcommand{\Acal}{\mathcal{A}}
\newcommand{\Scal}{\mathcal{S}}
\newcommand{\Hcal}{\mathcal{H}}
\newcommand{\Eop}{\mathbb{E}}
\newcommand{\Tset}{T^{\mathrm{ep}}}

\title{\textbf{Decision Modeling for Predict--Detect--Simulate Logistics:\\
A Pre-Factory-Tour Framework for DENSO Factory Hacks D2}\\[6pt]
\large Constraints, Objectives, Uncertainty, and Simulation-Based Recommendation}
\author{Ph\d{a}m Quang Minh\\[2pt]
\small Decision/Optimization/Mathematical Modeling role, DENSO Factory Hacks 2026 D2}
\date{September 2026}

\begin{document}
\maketitle

\begin{abstract}
\noindent This report is the normative specification of the decision layer of
the D2 system (``Simulate \& Forecast Logistics, Recommend Actions'') as it can
be defined before any factory-specific evidence exists. The recommendation
problem is stated over a candidate action set filtered by a three-valued
(satisfied/violated/unknown) hard-constraint logic, evaluated through a
discrete-event simulator under a shared scenario--seed bank, and aggregated by
a preference rule that remains a human decision. Six candidate decision
families, a four-class constraint taxonomy, and the simulator--optimizer
ownership boundary are fixed; the actual family, objective, and cadence wait
for the Factory Tour (2026-09-11). All factory-specific quantities are labeled
synthetic or unknown, and no claim is made about DENSO's actual operations.
A C++20 core at \texttt{cpp/} implements the feasibility semantics as the
approved executable specification of this report.
\end{abstract}

\tableofcontents
\newpage

%======================================================================
\section{Introduction}\label{sec:intro}
%======================================================================

The D2 competition task requires a system that predicts future logistics
state, detects bottlenecks, simulates what-if scenarios, and recommends
feasible, useful actions. The conceptual pipeline is

\begin{center}
\fbox{\parbox{0.95\linewidth}{\centering
factory state / forecast $\rightarrow$ \textbf{forecast} $\rightarrow$
\textbf{simulation + bottleneck detection} $\rightarrow$ \textbf{candidate
actions} $\rightarrow$ \textbf{feasibility} $\rightarrow$ \textbf{simulated
KPI outcomes} $\rightarrow$ \textbf{ranking} $\rightarrow$
\textbf{recommendation}}}
\end{center}

This report specifies the decision layer only: given the current or forecast
state and simulator outcomes, \emph{which feasible action should be
recommended, and on what evidence}. Two roles are separated sharply. The
digital-twin/simulation layer (owned by Đào Văn Đức) models physical
evolution, event dynamics, KPI generation, and bottleneck behavior. The
decision layer (this report) owns candidate generation, feasibility modeling,
objective interpretation, and ranking. The decision layer never re-creates
simulator semantics, and the simulator never selects the recommended action
(\cref{sec:interface}).

The report is deliberately written \emph{before} the Factory Tour
(2026-09-11). Its purpose is to fix what can be modeled safely now ---
the feasibility semantics, the constraint taxonomy, the candidate decision
families, the objective structures, and the evaluation methodology --- and to
make explicit which facts must come from the plant before a final
formulation, objective, or algorithm can be chosen. Every factory-specific
quantity in this report is synthetic (an invented development value) or
unknown (missing information that blocks the claims depending on it); none is
a claim about DENSO's operations.

%======================================================================
\section{Scope, Assumptions, and Provenance}\label{sec:status}
%======================================================================

Two provenance vocabularies are used, and they are not interchangeable. The
\emph{data-parameter source types} are owned by the data layer and used
verbatim here:
\begin{equation}\label{eq:srcset}
\mathcal{S}_{\mathrm{src}}
=\{\texttt{synthetic},\ \texttt{observed},\ \texttt{estimated},\
\texttt{assumed},\ \texttt{mentor\_confirmed},\ \texttt{system\_export}\}.
\end{equation}
Separately, this report attaches \emph{claim statuses} to its own modeling
statements (\cref{tab:status}). Each parameter records a tuple
$\langle \text{value},\ \text{unit},\ \text{source},\ \text{confidence},\
\text{validation status} \rangle$; a parameter whose value, unit, or source is
not established is unknown and must not silently become a numerical
assumption.

\begin{table}[htbp]
\centering\small
\caption{Claim statuses and permitted use. Data-parameter provenance uses
the data-layer enum \eqref{eq:srcset}, not this table.}
\label{tab:status}
\begin{tabular}{@{}p{3.1cm}p{6.9cm}p{4.2cm}@{}}
\toprule
Status & Meaning & Can support a DENSO claim? \\
\midrule
GENERIC & Structural property valid across factories (capacities are finite)
& No --- structure only \\
SYNTHETIC & Invented value for development/testing (e.g., 2 transporters)
& No --- development only \\
CANDIDATE & Proposed structure awaiting Factory-Tour validation & No ---
design only \\
OBSERVED & Measured during the tour with provenance & Provisionally, with
sample basis \\
MENTOR-CONFIRMED & Validated by a DENSO mentor & Yes \\
UNKNOWN & Required information not yet available & No --- blocks dependent
claims \\
\bottomrule
\end{tabular}
\end{table}

The central discipline of this report is the unknown-information rule: missing
information about a required hard constraint yields status unknown, never
implicit feasibility,
\begin{equation}\label{eq:unknownrule}
\text{data missing for constraint } i
\;\Longrightarrow\;
\phi_i(a,s_t)=\texttt{unknown},
\qquad
a \notin \Acal_F(s_t),
\end{equation}
i.e., the affected action is blocked and flagged for mentor validation rather
than passed. This is the opposite of the common failure mode in which absent
constraints silently make every action feasible.

%======================================================================
\section{Problem Definition}\label{sec:problem}
%======================================================================

\subsection{States, Actions, and the Catalog}\label{sec:actions}

Let $s_t$ denote the decision-relevant state at epoch $t$, assembled from
simulation outputs (queue lengths, WIP, utilization, bottleneck label) and
the forecast trajectory. The decision layer proposes actions from a
state-dependent set $\Acal(s_t)$, drawn from an \emph{action catalog} whose
entries carry provenance (source type, confidence, validation status, mentor
flag) and are, before the tour, synthetic development artifacts. The catalog
contains the explicit no-action alternative $a_0$; ``no feasible
improvement'' is a valid, reportable outcome of every decision epoch, not a
failure.

Six generic action families are CANDIDATE structures for the catalog ---
resource reallocation, dispatch-priority change, route change,
buffer/capacity adjustment, replenishment-interval change, and temporary
capacity addition --- alongside the explicit no-op ($a_0$ itself). These
\emph{action families} (the typed payload shapes realized in the C++ core)
must not be confused with the \emph{decision families} of
\cref{sec:families}, which group whole decision problems. Concrete action
semantics (which resource, which route, which quantity) are unknown until
the tour; the C++ core fixes only the typed payload shapes, not any DENSO
instances.

\subsection{Candidate Decision Families}\label{sec:families}

The actual DENSO decision family is unknown before the tour. Six families
are formulated as candidates (\cref{tab:families}); structural taxonomies of
automotive intralogistics support the family menu
\cite{boysen2015part,leanh2006agv,fragapane2021amr}, but none of them is
evidence about DENSO's plant.

\begin{table}[htbp]
\centering\small
\caption{Candidate decision families. All remain CANDIDATE until the
Factory Tour; the priors on the first two rows are inferences from our own
synthetic 2-action development catalog, not evidence about DENSO.}
\label{tab:families}
\begin{tabular}{@{}p{2.6cm}p{2.6cm}p{3.4cm}p{3.0cm}p{2.6cm}@{}}
\toprule
Family & Decision variables & Data needed & First method & Tour trigger \\
\midrule
Resource reallocation & $x_{rj} \in \{0,1\}$ & counts, ceilings
$\bar{c}_r$, compatibility & assignment / MILP & resources re-assignable at
operational cadence \\
Dispatching & $\pi \in \Pi$ & current rule, allowed variants & catalog
enumeration & dispatch rules adjustable \\
Routing & $y_{r,e,t} \in \{0,1\}$ & route graph, travel times & exact + local
search & routing freedom confirmed \\
Replenishment & $\rho_j, \beta_j$ & policy bounds, demand link & policy
enumeration / MILP & line feeding confirmed \\
Buffer / capacity & $\Delta C_n$, $u_r$ & limits, authorization, cost & small
integer set & buffer adjustment executable \\
Integrated prod.--transport & coupled scheduling variables & schedule
interface, coupling & exact small / matheuristic & strong coupling (S7) \\
\bottomrule
\end{tabular}
\end{table}

Resource reallocation and dispatching carry a modest catalog-informed prior:
one of the two synthetic development actions is a resource-count change, and
dispatch-rule selection is the other lever the current pipeline can express.
For reallocation, the minimal structure is an assignment problem: with
$x_{rj} = 1$ iff resource $r$ serves job $j$, resource exclusivity and job
coverage read
\begin{equation}\label{eq:exclusivity}
\sum_{j \in J} x_{rj} \;\le\; 1 \quad \forall\, r \in R,
\qquad\qquad
\sum_{r \in R} x_{rj} \;=\; 1 \quad \forall\, j \in J,
\end{equation}
which becomes a generalized assignment problem once capacity or travel-time
costs enter. For dispatching, the rule $\pi$ is chosen from a small catalog
$\Pi$ (FIFO, earliest-deadline-first, bottleneck-first); at single-digit
catalog sizes, exhaustive evaluation of all rules is both optimal within the
catalog and fully explainable, so sophisticated heuristics add nothing.

The routing family deserves one caution: the AGV/AMR literature
\cite{leanh2006agv,fragapane2021amr} is an analogue, not proof that D2 is an
AGV routing problem --- if the tour confirms fixed routes and tugger-train
operations, the family collapses into replenishment. The integrated
production--transport family is conditional on tour question S7; published
approaches at scale are reported to use matheuristics or decomposition, but
those pointers are unverified in our survey and no claim here is
source-supported. It is substantially harder, because production and
transport variables must be coupled over time, and should be selected only
if the coupling is real.

\subsection{Simulator Interface}\label{sec:interface}

The simulator is modeled as a map
\begin{equation}\label{eq:simulator}
Y : (s_t,\ a,\ \xi_s,\ \omega) \;\longmapsto\; \mathbf{v} \in \mathbb{R}^{K},
\end{equation}
returning the KPI vector realized under action $a$, scenario $\xi_s$, and
replication seed $\omega$. The decision layer consumes
\texttt{SimulationResult} objects only through this contract --- with
per-run metadata (scenario ID, seed, KPI units, action ID) --- and all
physical semantics stay behind it.

\begin{center}
\small
\begin{tabular}{@{}p{7.2cm}p{7.2cm}@{}}
\toprule
\textbf{Decision layer owns} & \textbf{Simulation layer owns} \\
\midrule
Candidate generation ($\Acal(s_t)$) & Physical evolution of the system \\
Feasibility filtering ($\phi_i$, verdicts) & Event dynamics, queues,
capacities \\
Objective interpretation ($F$) & KPI value generation (definitions per Red
decision) \\
Ranking and recommendation & Bottleneck detection behavior \\
Provenance labels on recommendations & Seeds, streams, replications \\
\bottomrule
\end{tabular}
\end{center}

The decision layer must not re-create simulator semantics, and the
simulation layer must not silently select the final recommended action. This
split mirrors the constraint-split architecture verified in recent
simulation-optimization practice \cite{ghaedyheidary2024simopt}: the
optimizer enforces deterministic feasibility and proposes partial decisions;
the simulator completes them and evaluates outcomes. Action masking under a
formal system model is the same principle \cite{lassoued2025fms}.

%======================================================================
\section{Mathematical Model}\label{sec:model}
%======================================================================

\cref{tab:notation} fixes the notation. Sets $R, N, E, J$ are generic
structures; their instances (which resources exist, which routes are
admissible) are unknown until the tour, and no set is instantiated before its
family is selected.

\begin{table}[htbp]
\centering\small
\caption{Notation. Only symbols used in this report are listed.}
\label{tab:notation}
\begin{tabular}{@{}lp{12.5cm}@{}}
\toprule
Symbol & Meaning \\
\midrule
$N, R, J, E$ & Nodes (stations, buffers, supermarkets); resources
(transporters, AGVs, operators); jobs; admissible routes \\
$A$, $\Acal(s_t)$ & Action catalog; actions applicable in state $s_t$ \\
$\Tset$, $t$, $\Delta$ & Decision epochs, epoch index, epoch length; horizon
$H = |\Tset|$ \\
$\Scal$, $\xi_s$ & Uncertainty scenarios; joint trajectory (forecast path +
parameter set) \\
$\omega$ & Simulation replication seed \\
$K$ & KPI dimensions \\
$\Hcal$, $g_i$, $\phi_i$ & Hard constraints; mathematical form $g_i \le 0$;
three-valued status function \\
$s_t$ & System state at epoch $t$ (queues, WIP, utilization, bottleneck
label, forecast trajectory) \\
$c_r$, $\delta_r$, $\bar{c}_r$ & Current count of resource $r$; count change
proposed by an action; pool ceiling \\
$C_n$, $q_n(t)$ & Capacity of node $n$; queue/WIP content at node $n$ \\
$\tau_e$, $p_j$ & Travel time on route $e$; processing time of job $j$ \\
$c_a$, $B$ & Implementation cost of action $a$; per-epoch budget \\
$x_{rj}$, $\pi$, $y_{r,e,t}$, $\rho_j$, $\beta_j$, $u_r$, $\Delta C_n$ &
Family decision variables (\cref{tab:families}) \\
$Y$, $\mathbf{y}(a)$, $\Delta_s(a)$ & Simulator map \eqref{eq:simulator};
KPI outcome vector; KPI delta vs.\ no-action under scenario $s$ \\
$F$, $a^{*}$ & Preference functional; recommended action \\
\bottomrule
\end{tabular}
\end{table}

\subsection{Feasibility Semantics}\label{sec:hardsoft}

Each hard constraint $i \in \Hcal$ carries a status function
\begin{equation}\label{eq:phistatus}
\phi_i(a, s_t) \in \{\texttt{satisfied},\ \texttt{violated},\ \texttt{unknown}\},
\end{equation}
mapping the mathematical form $g_i(a,s_t) \le 0$ to \texttt{satisfied} when
$g_i$ is verified evaluable with $g_i(a,s_t) \le 0$, to \texttt{violated} when
verified evaluable with $g_i(a,s_t) > 0$, and to \texttt{unknown} when the
data required to evaluate $g_i$ is missing. The action-level verdict composes
these values with three-valued logic:
\begin{equation}\label{eq:verdict}
\text{verdict}(a)
=
\begin{cases}
\texttt{INFEASIBLE} & \exists\, i \in \Hcal : \phi_i(a,s_t) = \texttt{violated},\\[2pt]
\texttt{UNKNOWN} & \phi_i(a,s_t) \ne \texttt{violated}\ \forall i \in \Hcal
\text{ and } \exists\, i \in \Hcal : \phi_i(a,s_t) = \texttt{unknown},\\[2pt]
\texttt{FEASIBLE} & \phi_i(a,s_t) = \texttt{satisfied}\ \forall\, i \in \Hcal.
\end{cases}
\end{equation}
Violation dominates; unverifiable constraints surface as unknown rather than
being silently passed; only fully verified satisfaction yields feasibility.
The feasible set is therefore
\begin{equation}\label{eq:feasibleset}
\Acal_F(s_t)
=
\left\{
a \in \Acal(s_t)
\;:\;
\phi_i(a, s_t) = \texttt{satisfied} \quad \forall\, i \in \Hcal
\right\}.
\end{equation}

The no-action alternative $a_0$ is feasible \emph{by construction} in the
following evidence-aware sense. Change-dependent constraints take their
$\delta = 0$ branch, so unknown data about proposed changes (pool ceilings,
budgets, lead times) never blocks the status quo. State-dependent
constraints, however, evaluate $a_0$ for real: if verified state evidence
shows the status quo already violates such a constraint (an evidenced buffer
overflow), the violation is reported for $a_0$ exactly as
\cref{eq:phistatus} requires --- satisfaction is never fabricated from
missing change data, and known violations are never suppressed.

Soft constraints do not participate in the verdict of \eqref{eq:verdict}. A
constraint may be declared soft by a named owner; its violation then appears
as an explicit flag in the recommendation reason, never as a silent score
shift, and hard-constraint violations are never expressed as objective
penalties. One metamorphic property is used throughout the pre-tour test
suite: adding hard constraints (in particular, tightening one) cannot enlarge
the feasible set,
\begin{equation}\label{eq:monotone}
\Hcal' \supseteq \Hcal
\;\Longrightarrow\;
\Acal_F(s_t; \Hcal') \subseteq \Acal_F(s_t; \Hcal).
\end{equation}

With feasibility fixed, the recommendation problem is
\begin{equation}\label{eq:generic-problem}
a^{*}
\in
\arg\min_{a \in \Acal_F(s_t)}
F\!\left(
\{\Delta_s(a)\}_{s \in \Scal}
\right),
\qquad
\Delta_s(a) := Y(s_t, a, \xi_s, \omega) - Y(s_t, a_0, \xi_s, \omega),
\end{equation}
where $F$ is a preference functional over scenario-indexed KPI deltas
(\cref{sec:objectives}). Deliberately, \eqref{eq:generic-problem} does not
commit to a scalar objective, and selecting $F$ (weights, orders,
thresholds) is a human decision, not an agent inference.

\subsection{Constraint Classes}\label{sec:constraints}

Constraints fall into four classes: \emph{physical} (the modeled system
itself; semantics owned by the simulation layer), \emph{operational} (how
the plant is run), \emph{business} (budget, labor, service commitments), and
\emph{decision-policy} (self-governance rules of the decision layer itself,
e.g.\ ``no action is labeled feasible with unvalidated provenance''). The full
pre-tour registry with identifiers, forms, and data needs is in
\cref{sec:registry}; two structural aspects belong here.

First, every constraint carries a \emph{check layer}: \emph{static}
(evaluable from configuration before simulation), \emph{simulator-enforced}
(an invariant the simulator must maintain --- resource exclusivity, buffer
bounds, flow conservation, availability windows --- where the decision layer
can only report violations observed in simulation results), or \emph{dynamic}
(checked on simulation outputs). For a simulator-enforced constraint, the
decision layer does not evaluate $\phi_i$ from configuration; it remains
\texttt{unknown} unless simulation evidence is supplied.

Second, static pre-checks are admissible in one direction only. When the
\emph{proposed change} alone makes satisfaction impossible --- an
\texttt{ADJUST\_BUFFER} action requesting a negative target level, or a
proposed resource count exceeding a known pool ceiling,
$0 \le c_r + \delta_r \le \bar{c}_r$ --- the decision layer may statically
detect the violation at proposal time. Such checks turn \texttt{unknown} into
\texttt{violated}, never into \texttt{satisfied}: verifying that a
simulator-enforced constraint \emph{held} requires simulation evidence.
C-OPS-007 (action lead time) similarly admits a static deadline check
whenever the epoch carries a known response deadline.

\subsection{Objective Structures}\label{sec:objectives}

The KPI set is provisional; each KPI must eventually carry definition, unit,
aggregation window, improvement direction, and DENSO priority, none of which
is confirmed. The candidate KPI vector (deltas relative to no-action) is
\begin{equation}\label{eq:kpivec}
\mathbf{y}(a)
=
\bigl[
\Delta\,\mathrm{throughput},\;
\Delta\,\mathrm{lead\ time},\;
\Delta\,\mathrm{lateness},\;
\Delta\,\mathrm{WIP},\;
\Delta\,\mathrm{cost},\;
\mathrm{risk}
\bigr],
\end{equation}
where risk is a placeholder for scenario-spread summaries. Queue length,
service level, and stability may join the vector after tour confirmation;
utilization is deliberately excluded, because high utilization at a
constraint is efficient while at non-constraints it signals fragility --- its
direction is context-dependent and its priority is a DENSO decision.

Two preference structures are the primary candidates. The constrained single
objective,
\begin{equation}\label{eq:constrained}
\min_{a \in \Acal_F(s_t)} c_a
\quad\text{s.t.}\quad
L(a) \le L_{\max},
\end{equation}
is appropriate when one operational limit dominates and the threshold
$L_{\max}$ is defensible; the threshold itself is a human decision. The
lexicographic objective orders KPIs by strict priority $k_1 \succ k_2 \succ
\dots$: action $a$ dominates $b$ if
\begin{equation}\label{eq:lex}
f_{k_1}(a) < f_{k_1}(b)
\quad\text{or}\quad
\bigl(f_{k_1}(a) = f_{k_1}(b)\ \text{within tolerance }\varepsilon_1
\text{ and } (f_{k_2}(a),\dots) \succ_{\mathrm{lex}} (f_{k_2}(b),\dots)\bigr).
\end{equation}
Tolerances $\varepsilon_k$ must be explicit; without them, negligible primary
gains can justify large secondary losses. Lexicographic methods are standard
preference-modeling tools \cite{miettinen1998nonlinear}; whether DENSO
priorities are genuinely strict is a tour question.

\begin{table}[htbp]
\centering\small
\caption{Alternative preference structures, kept as candidates. None is
adopted before the tour.}
\label{tab:objectives}
\begin{tabular}{@{}p{2.9cm}p{4.6cm}p{3.6cm}p{3.6cm}@{}}
\toprule
Structure & Form & Appropriate when & Caution \\
\midrule
Epsilon-constraint & $\min f_1$ s.t.\ $f_k \le \epsilon_k$ & thresholds are
operationally interpretable & thresholds are policy decisions; infeasible
combinations must be reported \cite{miettinen1998nonlinear} \\
Weighted normalized & $F = \sum_k w_k \hat{f}_k$ & preferences are genuinely
compensatory and weights are elicited & incompatible units; scale sensitivity;
unstable under $\pm 20\%$ weight perturbation \\
Pareto screening & dominance over $\mathbf{y}(a)$ & first screen before
weights exist & a frontier is a shortlist, not a recommendation; noise
corrupts dominance at small sample sizes \\
\bottomrule
\end{tabular}
\end{table}

Evolutionary frontier search such as NSGA-II \cite{deb2002nsga2} is justified
only when enumeration is impractical and noise is controlled; with tens of
candidates and noisy simulation outputs, enumeration with replication
discipline is usually superior. Risk-aware aggregation is presented but
deferred: given scenario losses $C_s(a)$, the candidates are the expectation
$\Eop[C(a)] = \sum_s p_s C_s(a)$, the worst case $\max_s C_s(a)$, and the
CVaR representation \cite{rockafellar2000cvar}
\begin{equation}\label{eq:cvar}
\operatorname{CVaR}_{\alpha}(C(a)) =
\min_{t}\Bigl[\, t + \tfrac{1}{1-\alpha}\, \Eop\bigl[(C(a) - t)_+\bigr] \Bigr].
\end{equation}
Escalating a risk measure from a reported column to a decisive rule requires
(i) asymmetric downside consequences, (ii) service-level violations that
matter, and (iii) sufficient tail data --- empirical CVaR at level $\alpha$
from $N$ samples is driven by only $\approx N(1-\alpha)$ tail observations.
None of the three holds today: the pipeline is deterministic. The pre-tour
policy is report-first --- rank on expected KPI, report observed worst case
and (when tail support suffices) CVaR as secondary columns.

Cost semantics deserve one explicit rule: \emph{known-zero cost and unknown
cost are distinct states}. An action whose cost has not been recorded is not
a free action; ranking rules may consume cost only when its value and unit
are known (mentor-confirmed or synthetic by design). If absolute costs stay
undisclosed, an explicitly labeled ordinal ranking (cheap/medium/expensive)
is acceptable, and the objective must then use forms that do not need a
common unit --- lexicographic or epsilon-constraint, never a fabricated
currency sum.

\subsection{Uncertainty and Scenario Evaluation}\label{sec:uncertainty}

Three uncertainty sources are kept distinct: (i) \emph{forecast uncertainty}
--- the forecast contract provides quantile paths (p10/p50/p90; current code
consumes only p50), and quantiles are not scenarios with probabilities;
(ii) \emph{process randomness} --- service, travel, and breakdown variability
inside the simulator, which exists only once the DES is stochastic; and
(iii) \emph{parameter uncertainty} --- tour-measured values may be wrong or
drift, handled by labeled ranges and sensitivity analysis. Forecast residual
uncertainty and simulator variability can model the same mechanism twice;
scenario construction must state which variability lives where. Constructing
valid joint scenario trajectories is a forecast-layer task: trajectories must
preserve serial correlation across the horizon and cross-variable dependence
(independent per-period sampling from marginals creates impossible
trajectories), and any scenario weights require a sourced justification.

Actions are compared on a \emph{shared scenario--seed bank}: the identical
scenario set and seed set across actions wherever possible. The same integer
seed does not by itself guarantee valid pairing once actions alter event
paths, so pairing requires either a common-random-numbers mechanism (a
simulator-design question owned by the simulation layer) or adequate
independent replications with explicit unpaired analysis. Reporting uses
mean, variance, quantiles, and confidence intervals \cite{alexopoulos2006output}.
Quantile paths are used as labeled stress tests --- a candidate is evaluated
under each path and rank stability across paths is reported --- never
averaged into a single number, which would manufacture a probability
interpretation the forecast never provided. Sample-average approximation
formalizes Monte Carlo evaluation of expectations when scenario sets are
sampled \cite{amaran2016simulation,kleywegt2002saa}; robust-optimization and
scenario-approach guarantees are explicitly not transferred to our black-box
DES setting without their own assumptions
\cite{bertsimas2004price,campi2008scenario}.

%======================================================================
\section{Solution and Evaluation Strategy}\label{sec:strategy}
%======================================================================

\subsection{Baselines}\label{sec:baselines}

\begin{table}[htbp]
\centering\small
\caption{Decision-method ladder. Every stronger method must beat the
mandatory baselines under the same scenario--seed bank before it is
considered.}
\label{tab:baselines}
\begin{tabular}{@{}p{4.2cm}p{6.4cm}p{3.8cm}@{}}
\toprule
Method & Role & Placement rationale \\
\midrule
Baseline 0: no action & Null recommendation; every claim is a delta vs.\ this
& Mandatory reference \\
Baseline 1: rule-based heuristic & ``if bottleneck $=$ transport, add
transport capacity''-style rules & Mandatory; simple, explainable \\
Baseline 2: greedy enumeration & Evaluate all feasible single actions; rank
by approved rule & Likely first real optimizer; optimal within the catalog \\
\addlinespace
Exact reference & Exhaustive enumeration on tiny instances & Correctness
oracle (golden cases) \\
Structured exact & MILP / CP-SAT / min-cost flow / assignment & Only if the
family has algebraic structure \\
Conditional larger search & Local search, tabu, ILS, ALNS & Only when
$|\Acal| \gg$ budget and the family survives selection \\
Deferred & NSGA-II, RL, matheuristics, DRO & Not justified pre-tour (catalog
small; assumptions unverified) \\
\bottomrule
\end{tabular}
\end{table}

\paragraph{Selection rule (pre-tour default).} The Baseline 2 pipeline
recommends a candidate only when it is \emph{strictly} preferred over the
no-action alternative $a_0$ under the declared policy: a tie with $a_0$
keeps the status quo, because a zero-improvement change with unknown cost
and lead time is not a recommendation. $a_0$ itself is evaluated through
the same gate; when evidence shows the status quo violates a hard
constraint, $a_0$ is infeasible --- it is then neither simulated nor
compared against, and the best feasible candidate is recommended while the
$a_0$ entry exposes the violation. When nothing is feasible, the outcome
is an explicit no-feasible-option report, never a fabricated
recommendation. This rule is a conservative pre-tour default and is a
Factory-Tour confirmation item, not a settled preference.

\paragraph{Post-tour logistics specialization.}
The executable logistics slice keeps the generic gate unchanged and applies
the same three-valued semantics in two explicit phases. Its typed
single-action catalog is $a_0$, prioritize a transport task, assign a task to
an AMR, reassign a task between AMRs, and select an alternate route. Candidate
order is canonical by family and typed identifier, independent of container
insertion order. Pre-simulation admission verifies only snapshot-evaluable
hard constraints; a hard violation or missing required static evidence is not
simulated. For an admitted action $a$, post-simulation evidence is scoped by
the tuple
\[
(\text{snapshot id},\text{snapshot version},a,\omega,r,
  \text{model version}).
\]
An identity mismatch is an input error, not a feasibility observation.
Missing required post-simulation evidence yields \textsc{Unknown}, a verified
hard violation yields \textsc{Infeasible}, and only final
\textsc{Feasible} outcomes enter comparison. Battery thresholds, deadline
lead times, compatibility evidence, dynamic limits, and the KPI preference
remain explicit caller-owned inputs; this specialization encodes no
factory-specific value or objective weight. The no-action alternative passes
through the same simulation and post-gate path, and a policy tie retains the
status quo. If final $a_0$ evidence is \textsc{Unknown}, the pipeline reports
a distinct baseline-evidence-unknown status and makes no recommendation; this
is not the same claim as ``no feasible option.'' The built-in post-tour test
policy ignores optional cost and forms a total preorder over finite
throughput, lead-time, lateness, and WIP outcomes. This is synthetic test
machinery, not a DENSO KPI hierarchy.

For small discrete candidate sets, enumeration with statistically careful
comparison is the methodologically indicated choice
\cite{amaran2016simulation,kleywegt2002saa}; recent results show that for
small recurring re-decisions even exact solutions can be milliseconds away
\cite{adamo2024realtime}, so search sophistication is not the bottleneck at
our scale. Ranking-and-selection machinery for many noisy alternatives is
deferred until needed \cite{wang2024sequential}.

One capability gap is acknowledged plainly: the current repository simulator
accepts only a configuration and a forecast --- it cannot accept an action
--- so simulation-based ranking is blocked until the simulation owner
delivers an action-capable simulator. Until then, the implementable
decision-layer work is the feasibility gate, the provenance-labeled action
catalog, and ranking over synthetic analytic outcomes, never fabricated
simulation results.

\subsection{Exact Small-Instance Oracle}\label{sec:oracle}

For small synthetic instances (2--3 resources, 3--5 actions), all feasible
single actions are enumerated and the exact optimum $f^{*}$ of the model
under evaluation is computed. A ranking method produces $f_h$; the
optimality gap is
\begin{equation}\label{eq:gap}
\mathrm{gap}(f_h, f^{*})
=
\begin{cases}
\dfrac{f_h - f^{*}}{|f^{*}|}, & f^{*} \ne 0,\\[8pt]
f_h - f^{*}, & f^{*} = 0 \quad \text{(zero-denominator guard)}.
\end{cases}
\end{equation}
This is a correctness benchmark on synthetic instances only --- never a
DENSO performance result. Its purposes are to verify that feasibility,
ranking, and preference-rule modules reproduce hand-computable optima, to
characterize runtime scaling honestly, and to serve as a differential-testing
oracle for any future optimized search kernel.

\subsection{Simulation Evaluation Budget}\label{sec:budget}

The evaluation cost per decision epoch is approximately
\begin{equation}\label{eq:evalcost}
\text{Cost} \;\propto\; |\Acal| \times |\Scal| \times R \times t_{\mathrm{sim}},
\end{equation}
with $R$ replications per (action, scenario) and $t_{\mathrm{sim}}$ runtime
per run. With a toy simulator and a small catalog the cost is negligible; it
becomes binding once the product reaches the order of hundreds to thousands
of full runs per epoch (e.g., $20 \times 10 \times 5 = 1000$), or when
$t_{\mathrm{sim}}$ grows with a real DES. The first response is not
optimization but the two-tier budget pattern verified in recent simulation
optimization \cite{ghaedyheidary2024simopt}: a small replication budget per
candidate during ranking ($R = 5$--$10$; design proposal, not a universal
constant), followed by a larger \emph{fresh-sample} validation of the
selected finalist ($\ge 10R$) that is never tuned on and used only for
confirmation. Sequential elimination \cite{wang2024sequential}, budget
allocation inside metaheuristics, common-random-numbers variance reduction,
memoization, and surrogates are deferred until a seeded DES and a measured
need exist.

\subsection{Algorithm Selection and Implementation Status}\label{sec:algo}

Algorithm selection follows problem identification, never precedes it.
Confirmed problem shape maps to a first method family: a small action
catalog maps to enumeration/greedy with paired evaluation; structured
assignment to MILP/CP-SAT/min-cost flow; routing and scheduling to an exact
baseline plus neighborhood search; a large black-box DES problem to
simulation optimization with a fixed budget and final validation; a dynamic
decision cadence to rolling-horizon or event-triggered re-solving. MILP,
CP-SAT, tabu, ILS, ALNS, genetic algorithms, and reinforcement learning are
not selected prematurely, and nothing native is justified at current problem
sizes.

Python remains the reference language for the pipeline (model, baselines,
orchestration, small exact instances); native kernels are a post-profiling
decision --- C++ when the solver ecosystem, combinatorial structure, or team
familiarity dominates, Rust when a larger custom parallel search engine with
concurrency-safety requirements is warranted. No native \emph{search kernel}
is planned in the competition window. A C++20 core at \texttt{cpp/}
implements the feasibility semantics of \cref{sec:hardsoft} and the catalog
as the approved executable specification of this report's pre-tour
semantics; that approval covers the executable-specification role, not a
choice of C++ for future search kernels. Its action-cost handling
deliberately differs from the Python shared contract --- the C++ core omits
cost entirely because its unit is unknown pre-tour, whereas the Python
dataclass defaults \texttt{estimated\_cost} to 0.0 --- and reconciling the
two is a shared-contract decision that must not be made silently by either
layer.

%======================================================================
\section{Data and Validation Requirements}\label{sec:data}
%======================================================================

\begin{center}
\small
\begin{longtable}{@{}p{3.4cm}p{3.4cm}p{1.4cm}p{2.4cm}p{3.2cm}@{}}
\caption{Compact data-requirement map.}\\
\toprule
Field & Used by & Required & Unit & Status / tour need \\
\midrule
\endfirsthead
\toprule
Field & Used by & Required & Unit & Status / tour need \\
\midrule
\endhead
Resource IDs & assignment family & yes & --- & UNKNOWN (structure only) \\
Route topology & routing, connectivity & yes & graph & UNKNOWN \\
Travel time $\tau_e$ & simulation, actions & yes & min & SYNTHETIC \\
Loading/unloading time & simulation & yes & min & UNKNOWN \\
Resource count $c_r$ & feasibility C-PHY-002 & yes & count & SYNTHETIC (2) \\
Pool ceiling $\bar{c}_r$ & ADD\_RESOURCE bound & yes & count & UNKNOWN \\
Buffer capacity $C_n$ & C-PHY-003 & yes & units & SYNTHETIC (20) \\
Queue/WIP $q_n$ & state $s_t$ & runtime & jobs & SYNTHETIC (toy simulator) \\
Demand forecast & scenario construction & yes & jobs/h & SYNTHETIC
(quantiles) \\
Replenishment policy & replenishment family & family & min/units & UNKNOWN \\
Deadline / service target & lateness KPI & objective & time & UNKNOWN \\
Shift schedule & C-OPS-002, epochs & yes & schedule & UNKNOWN \\
Action catalog & decision layer & yes & --- & SYNTHETIC (2 actions) \\
Action cost $c_a$ & C-BUS-001/003 & ranking & currency/ordinal & UNKNOWN \\
Action lead time & C-OPS-007 & honest ranking & min & UNKNOWN \\
Decision cadence & epoch design $\Tset$ & yes & per shift/day & UNKNOWN \\
KPI priorities & objective $F$ & yes & --- & UNKNOWN (Red decision) \\
\bottomrule
\end{longtable}
\end{center}

Validation proceeds in three layers before any optimization run: (A) schema
validation (fields, types, units, referential integrity); (B)
domain/physical checks --- non-negativity, capacity bounds, utilization in
$[0,1]$, interval ordering --- which are machine-checkable now; and (C)
statistical/temporal checks (distributions, gaps, drift, leakage) once
enough data exists. Impossible records are flagged, never silently repaired.
One reporting rule accompanies the KPI units once confirmed: a delta
smaller than the evaluation noise floor is not an improvement --- before an
action is called better than no-action, its paired delta must exceed the
dispersion of the paired differences (e.g., the confidence interval must
exclude zero) and matter operationally.

Synthetic development scenarios prove \emph{algorithmic} properties ---
correctness of feasibility logic, contract shapes, invariants, deterministic
behavior (same config + seed $\Rightarrow$ same result), and expected
directional responses --- and cannot prove real DENSO accuracy, ROI,
bottleneck identity, or optimality of any action. The six named scenarios
are \texttt{balanced\_flow} (no bottleneck; no action claims large
improvement), \texttt{known\_transport\_bottleneck} (transporter-constrained
state; a feasible \texttt{ADD\_RESOURCE} exists), \texttt{buffer\_overflow}
(queue at/over capacity; buffer-family and hard-constraint tests),
\texttt{demand\_spike} (p90 trajectory stress; capacity interaction),
\texttt{resource\_shortage} (pool ceiling makes \texttt{ADD\_RESOURCE}
infeasible via C-PHY-002), and \texttt{all\_actions\_infeasible} (the valid
output is an empty list with named violations, not a fake recommendation).
Metamorphic properties tested on them include \cref{eq:monotone}; increasing
bottleneck capacity must not decrease throughput in a deterministic
scenario; adding a dominated action must not change the winner; and missing
constraint data must block, not pass, the affected action.

%======================================================================
\section{Factory-Tour Decision Gates}\label{sec:tour}
%======================================================================

The Factory Tour (2026-09-11) resolves the core unknowns. Each question maps
to a specific part of the model:

\begin{center}
\small
\begin{tabular}{@{}p{6.6cm}p{7.6cm}@{}}
\toprule
Question & Affects \\
\midrule
M1. Actual logistics pain point? & Family selection
(\cref{sec:families}); everything downstream \\
M2. Which 2--3 actions are executable? & Action space $\Acal(s_t)$, action
catalog; feasibility data \\
M3. Which KPIs matter; current values? & Objective $F$
(\cref{sec:objectives}); improvement sizing \\
M4. What is forbidden? & Hard-constraint set $\Hcal$
(\cref{sec:hardsoft}, \cref{sec:registry}) \\
M5. Decision cadence / response deadline? & Static vs.\ rolling-horizon
formulation; epoch design $\Tset$; runtime budget \\
M6. One historical incident walkthrough & Ground-truth anchor for action
effects; validation \\
S1. Current manual decision rule? & Baseline 0/1 definition
(\cref{sec:baselines}) \\
S2. Cost structure (even ordinal)? & C-BUS constraints; objective cost term
(\cref{sec:objectives}) \\
S3. Resource flexibility / pool ceiling $\bar{c}_r$? & C-PHY-002 bound;
resource-family viability \\
S5. Action lead times / stability? & C-OPS-007/008 data; ranking honesty \\
S7. Production--transport coupling? & Whether the integrated family
(\cref{sec:families}) is live; simulator interface granularity \\
S8. Re-decision trigger in practice? & Rolling-horizon vs.\ event-triggered
vs.\ one-shot; epoch design (refines M5) \\
\multicolumn{2}{@{}p{14.2cm}@{}}{\small S4 (data availability/latency) and
S6 (demand pattern/forecast practice) feed the data and forecast layers,
not the decision model directly; they are captured in the tour checklist
but omitted here.} \\
\bottomrule
\end{tabular}
\end{center}

After the tour, each family is scored on relevance, data availability,
actionability, optimization structure, and simulator compatibility, with a
Go/No-Go gate per family. Bottleneck detection research
\cite{li2009bottleneck,lawrence1994shifting} motivates the cadence questions:
bottlenecks shift over time, so a one-shot formulation is unlikely to
survive contact with operations. All post-tour selection statuses are
provisional by construction until the tour answers arrive.

%======================================================================
\section{Limitations and Next Steps}\label{sec:limits}
%======================================================================

What this report fixes: the decision-layer architecture (generate
$\rightarrow$ gate $\rightarrow$ simulate $\rightarrow$ rank $\rightarrow$
recommend), the three-valued feasibility semantics
\eqref{eq:verdict}, the four-class constraint taxonomy with check layers,
the candidate objective structures with the preference rule left as a human
decision, the simulator--optimizer ownership boundary, and the evaluation
methodology (shared scenario--seed banks, paired deltas, two-tier
replication budget, fresh-sample confirmation).

What remains unknown: the actual decision family, the executable action
space, factory-specific hard constraints, the KPI hierarchy, decision
cadence, and real cost parameters --- all gated on the Factory Tour. The
report's honest limit is that it is a pre-tour artifact: it defines how a
recommendation must be produced and evaluated, not what DENSO's plant will
actually require.

The generic and post-tour logistics Baseline 0--2 pipelines are implemented
and tested in the C++ core
(`\texttt{cpp/}`): an explicit no-action baseline evaluated through the
same simulation boundary as every candidate, explicit rule-triggered and
catalog candidate generators, and greedy enumeration under an explicitly
synthetic test preference policy --- differential-tested against a tiny
exhaustive selection oracle. The logistics pipeline consumes action-scoped evidence
through a test seam pending an action-capable simulator from the simulation
owner. No final optimizer, heuristic, solver, or native search kernel is
attempted, and no synthetic test result is a DENSO performance claim.

%======================================================================
\appendix
\section{Constraint Registry}\label{sec:registry}
%======================================================================

This appendix lists the pre-tour constraint registry. Identifiers are stable
and used by the C++ core; forms are generic. Statuses are claim statuses
(\cref{tab:status}).

\subsection*{Physical constraints}

\begin{center}
\small
\begin{longtable}{@{}lp{4.0cm}p{3.4cm}p{1.8cm}p{2.3cm}@{}}
\caption{Physical constraints. Simulation semantics owned by the simulation
layer.} \label{tab:physical}\\
\toprule
ID & Constraint & Mathematical form & Layer & Status \\
\midrule
\endfirsthead
\toprule
ID & Constraint & Mathematical form & Layer & Status \\
\midrule
\endhead
C-PHY-001 & Resource exclusivity & $\sum_j x_{rj} \le 1\ \forall r$
\eqref{eq:exclusivity} & sim-enf. & GENERIC \\
C-PHY-002 & Resource-count bounds & $0 \le c_r + \delta_r \le \bar{c}_r$,
pool ceiling UNKNOWN & static & GENERIC / HARD \\
C-PHY-003 & Buffer bounds & $0 \le q_n(t) \le C_n\ \forall n,t$ &
sim-enf. & GENERIC / HARD \\
C-PHY-004 & Flow conservation & inflow $=$ outflow $+$ $\Delta$ storage on
every node & sim-enf. & GENERIC / HARD \\
C-PHY-005 & Route connectivity & transport only along $e \in E$ & static &
GENERIC \\
C-PHY-006 & Non-negativity & $\tau_e \ge 0$, $p_j \ge 0$, $q_n \ge 0$,
WIP $\ge 0$ & static & UNIVERSAL \\
C-PHY-007 & Availability windows & resource busy only within availability
calendar & sim-enf. & UNKNOWN \\
\bottomrule
\end{longtable}
\end{center}

\subsection*{Operational constraints}

\begin{center}
\small
\begin{longtable}{@{}lp{5.6cm}p{4.0cm}p{2.0cm}@{}}
\caption{Operational constraints (pre-tour candidates; most UNKNOWN).}\\
\toprule
ID & Constraint & Data needed & Status \\
\midrule
\endfirsthead
\toprule
ID & Constraint & Data needed & Status \\
\midrule
\endhead
C-OPS-001 & Only catalog actions recommendable & action catalog w/\
provenance & enforced now \\
C-OPS-002 & Shift windows (changes at boundaries only) & shift schedule &
UNKNOWN \\
C-OPS-003 & Dispatch restrictions (fixed priorities) & dispatch rules &
UNKNOWN \\
C-OPS-004 & Replenishment policy bounds & policy parameters & UNKNOWN \\
C-OPS-005 & Route/material restrictions & material--route compatibility &
UNKNOWN \\
C-OPS-006 & Safety: no unsafe congestion/interaction & layout, traffic
rules & UNKNOWN \\
C-OPS-007 & Action lead time (takes effect after implementation delay) &
per-action lead time & UNKNOWN \\
C-OPS-008 & Action cooldown / minimum stability period & cooldown policy &
UNKNOWN \\
\bottomrule
\end{longtable}
\end{center}

\noindent C-OPS-006 is hard by default: no action may create unsafe
congestion (e.g., AMR--pedestrian interaction --- synthetic illustration);
its data is unknown.

\subsection*{Business constraints}

\begin{center}
\small
\begin{tabular}{@{}lp{7.4cm}p{2.2cm}p{2.0cm}@{}}
\toprule
ID & Constraint & Hard/soft & Status \\
\midrule
C-BUS-001 & Budget: $\sum_{a \in \text{selected}} c_a \le B$ & hard if
budget real; else excluded & UNKNOWN \\
C-BUS-002 & Labor limits (overtime, headcount) & hard & UNKNOWN \\
C-BUS-003 & Implementation/switching cost counts & soft (flag) or objective
term & design proposal \\
C-BUS-004 & Maximum acceptable lateness/service level & hard if contractual
& UNKNOWN \\
\bottomrule
\end{tabular}
\end{center}

\subsection*{Decision-policy constraints (self-governance)}

\begin{center}
\small
\begin{tabular}{@{}lp{9.0cm}p{3.0cm}@{}}
\toprule
ID & Constraint & Status \\
\midrule
C-POL-001 & At most $N_{\mathrm{chg}}$ actions change per epoch
($\sum_a z_a \le N_{\mathrm{chg}}$) & $N_{\mathrm{chg}}$ UNKNOWN \\
C-POL-002 & Recommendation must include no-action comparison & enforced
(hard policy) \\
C-POL-003 & No action labeled feasible with unknown/unvalidated provenance &
enforced (hard policy) \\
C-POL-004 & Combined/bundled actions disabled until single-action effects
measured & enforced (hard policy, current) \\
C-POL-005 & Preference rule fixed only by named human decision & enforced
(hard policy) \\
C-POL-006 & No DENSO-confidential data in Git/outputs/demos & enforced (hard
policy) \\
\bottomrule
\end{tabular}
\end{center}

\bibliographystyle{plain}
\bibliography{references}

\end{document}
